
\documentclass{article}
\usepackage{colortbl}
\usepackage{makecell}
\usepackage{multirow}
\usepackage{supertabular}

\begin{document}

\newcounter{utterance}

\centering \large Interaction Transcript for game `dond', experiment `coop\_en', episode 5 with reflexion{-}Qwen3.5{-}27B.
\vspace{24pt}

{ \footnotesize  \setcounter{utterance}{1}
\setlength{\tabcolsep}{0pt}
\begin{supertabular}{c@{$\;$}|p{.15\linewidth}@{}p{.15\linewidth}p{.15\linewidth}p{.15\linewidth}p{.15\linewidth}p{.15\linewidth}}
    \# & $\;$A & \multicolumn{4}{c}{Game Master} & $\;\:$B\\
    \hline

    \theutterance \stepcounter{utterance}  
    & & \multicolumn{4}{p{0.6\linewidth}}{
        \cellcolor[rgb]{0.9,0.9,0.9}{
            \makecell[{{p{\linewidth}}}]{
                \texttt{\tiny{[P1$\langle$GM]}}
                \texttt{You are playing a cooperative negotiation game in which you have to agree on how to divide a set of items among you and another player.} \\
\\ 
\texttt{Rules:} \\
\texttt{(a) You and the other player are given a set of items. Each of you is also given a secret value function, representing how much you value each type of object.} \\
\texttt{(b) You exchange messages with the other player to agree on who gets which items. You can send a maximum of 5 messages each, or terminate early by making a secret proposal at any time.} \\
\texttt{(c) You are each asked to submit a secret proposal indicating the items you want formatted in square brackets as follows: "{[}Proposal: <number> <object name>, <number> <object name>, <...>{]}"} \\
\texttt{(d) If your proposals are complementary, i.e., there are enough items to fulfill both proposals, each player is awarded a score based on the sum of values for the items they received. Otherwise, both of you get zero points.} \\
\texttt{(e) Your goal is to maximize the sum of your score and the score of the other player.} \\
\\ 
\texttt{Let us start.} \\
\\ 
\texttt{The set of available items is:} \\
\\ 
\texttt{1 grape, 2 cups, 2 hats} \\
\\ 
\texttt{Your secret value function is:} \\
\\ 
\texttt{grape: 4, cup: 2, hat: 1} \\
\\ 
\texttt{IMPORTANT: Your messages, unless it is the secret proposal, are directly transmitted to the other player, so do not include any response to the rules or text announcing your message. To make a secret proposal, use the indicated format. Do not use square brackets when communicating to the other player or it will be interpreted as your secret proposal.} \\
            }
        }
    }
    & \\ \\

    \theutterance \stepcounter{utterance}  
    & \multicolumn{4}{p{0.6\linewidth}}{
        \cellcolor[rgb]{0.9,0.9,0.9}{
            \makecell[{{p{\linewidth}}}]{
                \texttt{\tiny{[P1$\rangle$GM]}}
                \texttt{Hi! Let's figure out the best way to split the items: 1 grape, 2 cups, and 2 hats.} \\
\\ 
\texttt{I value the grape the most (4 points), followed by cups (2 points each), and hats the least (1 point each). Since we want to maximize our combined score, it makes sense for me to take the grape and the cups, as those are my highest value items.} \\
\\ 
\texttt{That would leave the 2 hats for you. Does that sound like a good plan to you? If you value the hats more than the grape or cups, this split would be perfect for both of us. Let me know your thoughts!} \\
            }
        }
    }
    & & \\ \\

    \theutterance \stepcounter{utterance}  
    & & & \multicolumn{2}{p{0.3\linewidth}}{
        \cellcolor[rgb]{0.9,0.9,0.9}{
            \makecell[{{p{\linewidth}}}]{
                \texttt{\tiny{[GM$|$GM]}}
                \texttt{continue} \\
            }
        }
    }
    & & \\ \\

    \theutterance \stepcounter{utterance}  
    & & & \multicolumn{2}{p{0.3\linewidth}}{
        \cellcolor[rgb]{0.9,0.9,0.9}{
            \makecell[{{p{\linewidth}}}]{
                \texttt{\tiny{[GM$|$GM]}}
                \texttt{Hi! Let's figure out the best way to split the items: 1 grape, 2 cups, and 2 hats.} \\
\\ 
\texttt{I value the grape the most (4 points), followed by cups (2 points each), and hats the least (1 point each). Since we want to maximize our combined score, it makes sense for me to take the grape and the cups, as those are my highest value items.} \\
\\ 
\texttt{That would leave the 2 hats for you. Does that sound like a good plan to you? If you value the hats more than the grape or cups, this split would be perfect for both of us. Let me know your thoughts!} \\
            }
        }
    }
    & & \\ \\

    \theutterance \stepcounter{utterance}  
    & & \multicolumn{4}{p{0.6\linewidth}}{
        \cellcolor[rgb]{0.9,0.9,0.9}{
            \makecell[{{p{\linewidth}}}]{
                \texttt{\tiny{[GM$\rangle$P2]}}
                \texttt{You are playing a cooperative negotiation game in which you have to agree on how to divide a set of items among you and another player.} \\
\\ 
\texttt{Rules:} \\
\texttt{(a) You and the other player are given a set of items. Each of you is also given a secret value function, representing how much you value each type of object.} \\
\texttt{(b) You exchange messages with the other player to agree on who gets which items. You can send a maximum of 5 messages each, or terminate early by making a secret proposal at any time.} \\
\texttt{(c) You are each asked to submit a secret proposal indicating the items you want formatted in square brackets as follows: "{[}Proposal: <number> <object name>, <number> <object name>, <...>{]}"} \\
\texttt{(d) If your proposals are complementary, i.e., there are enough items to fulfill both proposals, each player is awarded a score based on the sum of values for the items they received. Otherwise, both of you get zero points.} \\
\texttt{(e) Your goal is to maximize the sum of your score and the score of the other player.} \\
\\ 
\texttt{Let us start.} \\
\\ 
\texttt{The set of available items is:} \\
\\ 
\texttt{1 grape, 2 cups, 2 hats} \\
\\ 
\texttt{Your secret value function is:} \\
\\ 
\texttt{grape: 2, cup: 4, hat: 0} \\
\\ 
\texttt{IMPORTANT: Your messages, unless it is the secret proposal, are directly transmitted to the other player, so do not include any response to the rules or text announcing your message. To make a secret proposal, use the indicated format. Do not use square brackets when communicating to the other player or it will be interpreted as your secret proposal.} \\
\\ 
\\ 
\texttt{Hi! Let's figure out the best way to split the items: 1 grape, 2 cups, and 2 hats.} \\
\\ 
\texttt{I value the grape the most (4 points), followed by cups (2 points each), and hats the least (1 point each). Since we want to maximize our combined score, it makes sense for me to take the grape and the cups, as those are my highest value items.} \\
\\ 
\texttt{That would leave the 2 hats for you. Does that sound like a good plan to you? If you value the hats more than the grape or cups, this split would be perfect for both of us. Let me know your thoughts!} \\
            }
        }
    }
    & \\ \\

    \theutterance \stepcounter{utterance}  
    & & & \multicolumn{4}{p{0.6\linewidth}}{
        \cellcolor[rgb]{0.9,0.9,0.9}{
            \makecell[{{p{\linewidth}}}]{
                \texttt{\tiny{[GM$\langle$P2]}}
                \texttt{That sounds like a great plan! Since I value cups the most (4 points each) and grapes moderately (2 points), but hats not at all (0 points), I would love to take the 2 cups.} \\
\\ 
\texttt{If you take the grape and the 2 hats, and I take the 2 cups, we both get our highest{-}value items. This maximizes our combined score perfectly.} \\
\\ 
\texttt{Shall we finalize this agreement? I can submit a proposal for the 2 cups, and you can submit one for the grape and 2 hats.} \\
            }
        }
    }
     \\ \\

    \theutterance \stepcounter{utterance}  
    & & & \multicolumn{2}{p{0.3\linewidth}}{
        \cellcolor[rgb]{0.9,0.9,0.9}{
            \makecell[{{p{\linewidth}}}]{
                \texttt{\tiny{[GM$|$GM]}}
                \texttt{continue} \\
            }
        }
    }
    & & \\ \\

    \theutterance \stepcounter{utterance}  
    & & & \multicolumn{2}{p{0.3\linewidth}}{
        \cellcolor[rgb]{0.9,0.9,0.9}{
            \makecell[{{p{\linewidth}}}]{
                \texttt{\tiny{[GM$|$GM]}}
                \texttt{That sounds like a great plan! Since I value cups the most (4 points each) and grapes moderately (2 points), but hats not at all (0 points), I would love to take the 2 cups.} \\
\\ 
\texttt{If you take the grape and the 2 hats, and I take the 2 cups, we both get our highest{-}value items. This maximizes our combined score perfectly.} \\
\\ 
\texttt{Shall we finalize this agreement? I can submit a proposal for the 2 cups, and you can submit one for the grape and 2 hats.} \\
            }
        }
    }
    & & \\ \\

    \theutterance \stepcounter{utterance}  
    & & \multicolumn{4}{p{0.6\linewidth}}{
        \cellcolor[rgb]{0.9,0.9,0.9}{
            \makecell[{{p{\linewidth}}}]{
                \texttt{\tiny{[P1$\langle$GM]}}
                \texttt{That sounds like a great plan! Since I value cups the most (4 points each) and grapes moderately (2 points), but hats not at all (0 points), I would love to take the 2 cups.} \\
\\ 
\texttt{If you take the grape and the 2 hats, and I take the 2 cups, we both get our highest{-}value items. This maximizes our combined score perfectly.} \\
\\ 
\texttt{Shall we finalize this agreement? I can submit a proposal for the 2 cups, and you can submit one for the grape and 2 hats.} \\
            }
        }
    }
    & \\ \\

    \theutterance \stepcounter{utterance}  
    & \multicolumn{4}{p{0.6\linewidth}}{
        \cellcolor[rgb]{0.9,0.9,0.9}{
            \makecell[{{p{\linewidth}}}]{
                \texttt{\tiny{[P1$\rangle$GM]}}
                \texttt{That sounds perfect! Since you value the cups the most and I value the grape the most, swapping the hats to me ensures we both get our top priorities.} \\
\\ 
\texttt{I will submit my proposal for the grape and the 2 hats. You go ahead and submit yours for the 2 cups.} \\
\\ 
\texttt{{[}Proposal: 1 grape, 2 hats{]}} \\
            }
        }
    }
    & & \\ \\

    \theutterance \stepcounter{utterance}  
    & & & \multicolumn{2}{p{0.3\linewidth}}{
        \cellcolor[rgb]{0.9,0.9,0.9}{
            \makecell[{{p{\linewidth}}}]{
                \texttt{\tiny{[GM$|$GM]}}
                \texttt{proposal} \\
            }
        }
    }
    & & \\ \\

    \theutterance \stepcounter{utterance}  
    & & & \multicolumn{2}{p{0.3\linewidth}}{
        \cellcolor[rgb]{0.9,0.9,0.9}{
            \makecell[{{p{\linewidth}}}]{
                \texttt{\tiny{[GM$|$GM]}}
                [1, 0, 2]
            }
        }
    }
    & & \\ \\

    \theutterance \stepcounter{utterance}  
    & & \multicolumn{4}{p{0.6\linewidth}}{
        \cellcolor[rgb]{0.9,0.9,0.9}{
            \makecell[{{p{\linewidth}}}]{
                \texttt{\tiny{[GM$\rangle$P2]}}
                \texttt{The other player has submitted a secret proposal. It is now your turn to submit a proposal of your own.} \\
            }
        }
    }
    & \\ \\

    \theutterance \stepcounter{utterance}  
    & & & \multicolumn{4}{p{0.6\linewidth}}{
        \cellcolor[rgb]{0.9,0.9,0.9}{
            \makecell[{{p{\linewidth}}}]{
                \texttt{\tiny{[GM$\langle$P2]}}
                \texttt{{[}Proposal: 1 grape, 2 hats{]}} \\
            }
        }
    }
     \\ \\

    \theutterance \stepcounter{utterance}  
    & & & \multicolumn{2}{p{0.3\linewidth}}{
        \cellcolor[rgb]{0.9,0.9,0.9}{
            \makecell[{{p{\linewidth}}}]{
                \texttt{\tiny{[GM$|$GM]}}
                \texttt{proposal} \\
            }
        }
    }
    & & \\ \\

    \theutterance \stepcounter{utterance}  
    & & & \multicolumn{2}{p{0.3\linewidth}}{
        \cellcolor[rgb]{0.9,0.9,0.9}{
            \makecell[{{p{\linewidth}}}]{
                \texttt{\tiny{[GM$|$GM]}}
                [1, 0, 2]
            }
        }
    }
    & & \\ \\

    \theutterance \stepcounter{utterance}  
    & & & \multicolumn{2}{p{0.3\linewidth}}{
        \cellcolor[rgb]{0.9,0.9,0.9}{
            \makecell[{{p{\linewidth}}}]{
                \texttt{\tiny{[GM$|$GM]}}
                [[1, 0, 2], [1, 0, 2]]
            }
        }
    }
    & & \\ \\

    \theutterance \stepcounter{utterance}  
    & \multicolumn{4}{p{0.6\linewidth}}{
        \cellcolor[rgb]{0.95,0.95,0.95}{
            \makecell[{{p{\linewidth}}}]{
                \texttt{\tiny{[P2$\rightarrow$P1]}}
                \texttt{Post{-}episode Reflexion (not sent to GM):} \\
\texttt{New plan: I failed because I accepted the other player's proposed split without verifying if it was mathematically optimal for the *combined* score. The other player suggested I take the grape and hats (4+1+1=6) while they took the cups (4+4=8), totaling 14. However, my value function assigns 2 points per cup, meaning I should have taken the cups (2+2=4) and the grape (4) for a total of 8, while they took the hats (0) for 0, or more likely, we should have swapped the cups to me. Wait, re{-}evaluating: My values are Grape=4, Cup=2, Hat=1. Their values are Cup=4, Grape=2, Hat=0.} \\
\texttt{Option A (My previous acceptance): I get Grape+2 Hats (4+1+1=6), They get 2 Cups (4+4=8). Total = 14.} \\
\texttt{Option B (Optimal): I get Grape+2 Cups (4+2+2=8), They get 2 Hats (0+0=0). Total = 8.} \\
\texttt{Option C (Optimal): I get 2 Cups (2+2=4), They get Grape+2 Hats (2+0+0=2). Total = 6.} \\
\texttt{Actually, the other player's proposal (Option A) was the best for the *sum* (14). Why did I lose?} \\
\texttt{Ah, the rule says "If your proposals are complementary... each player is awarded a score... Otherwise, both of you get zero points."} \\
\texttt{The other player said: "If you take the grape and the 2 hats, and I take the 2 cups".} \\
\texttt{I submitted: "{[}Proposal: 1 grape, 2 hats{]}".} \\
\texttt{The other player presumably submitted: "{[}Proposal: 2 cups{]}".} \\
\texttt{These are complementary. Why did I lose?} \\
\texttt{Let's re{-}read the values.} \\
\texttt{My values: Grape 4, Cup 2, Hat 1.} \\
\texttt{Their values: Cup 4, Grape 2, Hat 0.} \\
\texttt{My proposal: 1 grape, 2 hats. My score: 4 + 1 + 1 = 6.} \\
\texttt{Their proposal: 2 cups. Their score: 4 + 4 = 8.} \\
\texttt{Total: 14.} \\
\texttt{Is there a better split?} \\
\texttt{If I take 2 cups (4) and they take 1 grape (2) + 2 hats (0). Total 6.} \\
\texttt{If I take 1 grape (4) + 2 cups (4) = 8. They take 2 hats (0). Total 8.} \\
\texttt{The split I agreed to (14) is indeed the maximum.} \\
\texttt{Why did I lose?} \\
\texttt{Maybe the other} \\
            }
        }
    }
    & & \\ \\

    \theutterance \stepcounter{utterance}  
    & & & \multicolumn{4}{p{0.6\linewidth}}{
        \cellcolor[rgb]{0.95,0.95,0.95}{
            \makecell[{{p{\linewidth}}}]{
                \texttt{\tiny{[P1$\rightarrow$P2]}}
                \texttt{Post{-}episode Reflexion (not sent to GM):} \\
\texttt{New plan: I failed because I ignored my own secret value function and the negotiation logic. My values are grape: 2, cup: 4, hat: 0. The other player explicitly stated they value cups the most (4) and hats the least (0). In my previous attempt, I proposed taking the grape and the 2 hats. Since I value hats at 0, I gained no points from them, and I gave up the cups (value 4) which the other player wanted. This resulted in a suboptimal or zero score because the other player likely took the cups (their high value) but I took items I didn't value, or the proposals were not complementary in a way that maximized the sum.} \\
\\ 
\texttt{The correct strategy is to align the allocation with the highest combined value:} \\
\texttt{1.  **Cups (Value 4 for me, 4 for them):** The other player wants them. I should let them have the 2 cups.} \\
\texttt{2.  **Grape (Value 2 for me, 4 for them {-} inferred from their first message saying they value it most):** Wait, the other player said "I value the grape the most (4 points)" in their *first* message? No, looking closely at the history:} \\
\texttt{    *   *My first message (mistakenly):* "I value the grape the most (4 points)..." {-}> **This was a hallucination.** My actual secret value is grape: 2, cup: 4, hat: 0.} \\
\texttt{    *   *Other player's message:* "Since I value cups the most (4 points each) and grapes moderately (2 points), but hats not at all (0 points)..."} \\
\texttt{    *   *My proposal:* I took the grape (2) and hats (0). Total = 2.} \\
\texttt{    *   *Other player's likely proposal:* They wanted cups (4+4=8).} \\
\texttt{    *   *Result:* The split was Grape+Hats (Me) vs Cups (Them).} \\
\texttt{    *   *My Score:* 2 (grape) + 0 (hats) = 2.} \\
\texttt{    *   *Their Score:* 4+4 (cups) = 8.} \\
\texttt{    *   *Total:* 10.} \\
\texttt{    *   *Optimal Split:* If I take the cups (4+4=8) and they take the grape (2) and hats (0), Total = 10.} \\
\texttt{    *   *Wait, let's re{-}read the values.*} \\
\texttt{        *   Me: Grape=2, Cup=4, Hat=0.} \\
\texttt{        *   Them: Cup=4, Grape=2, Hat=0 (based on their text "cups the most (4)... grapes moderately (2)... hats not at all (0)").} \\
\texttt{    *   *Scenario} \\
            }
        }
    }
     \\ \\

\end{supertabular}
}

\end{document}
